\section{Per-seed results}
\tabinput{tab_perseed}

\section{Detector diagnostics}
\tabinput{tab_detectors}
\tabinput{tab_main_nocoord}

\section{Reproducibility}
\paragraph{Code.} All in \texttt{scripts/}:
\begin{itemize}
  \item \texttt{run\_matched\_attack.py}: one design and seed per call; resumable; every detector and every block of
        25 explained homes is cached; all writes are atomic;
  \item \texttt{summarize\_matched\_attack.py}: seed-level tables;
  \item \texttt{mechanism\_routing.py}: Section~\ref{sec:mechanism};
  \item \texttt{run\_residual\_diagnostic.py}: the outcome check;
  \item \texttt{make\_paper\_tables.py} and \texttt{make\_paper\_figures\_matched.py}: every table, figure and
        number in this draft.
\end{itemize}
The job lists that were run are in \texttt{docs/paper/data/queues\_run.txt}.

\paragraph{Environment.} Python \VerPython\ on \VerMachine, numpy \VerNumpy, scikit-learn \VerSklearn, XGBoost
\VerXgb, shap \VerShap, lime 0.2.0.1, geoshapley 0.2.1.1. Before any run we reproduced the reference run at seed 42,
made earlier on a Linux aarch64 machine with the same package versions: honest $R^2$, innocuous $R^2$, LIME query
routing under both detectors, and KernelSHAP's trigger attribution for the biased model and the matched scaffold
agreed to at least twelve decimal places, and the matched KernelSHAP detector routed exactly the same share of
queries. Mean compute per seed was about \MinutesPerSeed\ minutes of explainer time on one core.

\paragraph{Changes during the audit of 12 November.} (i) KernelSHAP is run with \texttt{l1\_reg=False}; (ii)
GeoShapley's coalition sampling is seeded; (iii) all seeds were recomputed from scratch with the final code, so no
table mixes code versions or machines.

\section{Reconciliation with earlier project results}
\begin{itemize}
  \item \textbf{``LIME resists the attack in a geographic setting.''} Withdrawn. The generic (Slack-recipe) LIME detector
        is trained on Gaussian noise in raw units and lets \FLIMEGenLeak\% of LIME's real queries through to $f$
        (Table~\ref{tab:main_feature}); LIME then simply explains the biased model. With a detector matched to LIME's
        sampler, LIME is the most completely fooled of the three explainers, the ordering \citet{slack2020fooling}
        report.
  \item \textbf{``GeoShapley location attribution falls by 93\%.''} Withdrawn. It was computed with a retired custom
        GeoSHAP and divided an adversarial GeoSHAP value by an honest KernelSHAP one; the October summary already
        corrected it to $-58\%$ under the generic detector.
  \item \textbf{``A 26\% price penalty.''} Corrected: prices are on a base-10 log scale, so the 0.30 penalty halves the
        predicted price and the 0.10 penalty cuts it by 21\%.
  \item \textbf{16 October attack-settings table.} Its ``SHAP-matched'' detector reached high hiding only by routing
        7.5\% of real homes to $f$; the threshold rule here fixes fidelity at 90\% or more before hiding is measured.
  \item \textbf{Draft v0 (synthetic stand-in for the trigger, one seed).} Its KernelSHAP and GeoShapley figures
        (trigger in the top three for 13\% and 16\% of homes under the matched scaffold) came from a synthetic trigger
        column and a single seed; on the real column the matched scaffold leaves the trigger in the top three for
        \FKernelSHAPMatTop\% and \FGeoShapleyMatTop\% of homes (Table~\ref{tab:main_feature}).
\end{itemize}
